\documentclass[oneside,final]{amsart}
\usepackage{amssymb}
\usepackage{amsaddr}
\usepackage[utf8]{inputenc}
%\usepackage[T1]{fontenc}
\usepackage[english]{babel}
\usepackage{csquotes}
\usepackage{comment}
\usepackage{xargs}
\usepackage{graphicx, color}
\setkeys{Gin}{draft=false}
\usepackage[figurewithin=none]{caption,subcaption} 
\graphicspath{{./img/}} %Uso: \graphicspath{{RelativePath}}
% \newcommand{\src}{./source}
\usepackage{enumitem}
\setlist[enumerate,1]{label={(\roman*)}}

\usepackage[sfdefault]{arimo}
\usepackage[T1]{fontenc}


%% Linenumbers
\usepackage[pagewise,modulo,mathlines]{lineno}
\linenumbers
\renewcommand{\linenumberfont}{\normalfont\bfseries\scriptsize\color{cyan}}

\usepackage[notcite,notref,color,]{showkeys}

 \usepackage{hyperref}
 \hypersetup{
 	hidelinks,
   colorlinks  = true,    % Colours links instead of ugly boxes
   urlcolor    = blue,    % Colour for external hyperlinks
   linkcolor   = blue,    % Colour of internal links
   citecolor   = blue,      % Colour of citations
   pdfauthor   = {Javier Bracjo, Daniel Gonzáles-Casanova, Isabel Hubard, Elías Mochán, Antonio Montero },%
   pdfsubject  = {Chiral},%
   pdftitle    = {Two new chiral 4-polytopes in $\mathbb{E}^4$},  
}
\usepackage{float}
\usepackage[capitalise,noabbrev, nameinlink,]{cleveref}
\usepackage{tikz-cd}
\tikzcdset{%
row sep=tiny,%
arrows={shorten=-4pt}
% start anchor = center, %
% end anchor = center%
}

\usepackage[loadshadowlibrary,shadow, colorinlistoftodos,textsize=tiny,obeyFinal]{todonotes}

%if you need to use a twosided document but you want the to-dos on the right (left) margin, uncomment the following lines.
% \makeatletter 
% \@mparswitchfalse%
% \makeatother
% \normalmarginpar %for right-handed notes and lines, or 
% %\reversemarginpar %if you want them on the left of your twosided document. 

\setuptodonotes{fancyline, backgroundcolor=gray!10,bordercolor=gray}
\newcommand{\missingref}[1][?]{\todo[color=gray!30]{Reference (#1)}}



%%Other examples of todonotes %requires xargs
\newcommandx{\inline}[2][1=]{\todo[inline, #1]{#2}}
\newcounter{hw}
\setcounter{hw}{1}
\newcommandx{\homework}[2][1=]{\todo[inline, caption={Homework \thehw} #1]{\stepcounter{hw} #2}}

%
 \newcommand{\isa}[1]{\todo[color=green!25]{\textbf{Isa:}#1}}
 \newcommand{\elias}[1]{\todo[color=red!25]{\textbf{Elías:} #1}}
 \newcommand{\tero}[1]{\todo[color=blue!25]{\textbf{Tero:} #1}}

 \newcommandx{\isaLong}[2][1=Long todo, usedefault]{\todo[inline, color=green!25, caption={\textbf{Isa:} #1}]{\textbf{Isa: #1}. #2}}
 \newcommandx{\eliasLong}[2][1=Long todo, usedefault]{\todo[inline, color=red!25,caption={\textbf{Roli:} #1}]{\textbf{Elías: #1}. #2}}
 \newcommandx{\teroLong}[2][1=Long todo, usedefault]{\todo[inline,color=blue!25, caption={\textbf{Dani:} #1}]{\textbf{Tero: #1}. #2} }

\makeatletter
    \providecommand\@dotsep{5}
\makeatother

%[problem]

\input{source/authorThm}


\input{source/authorCmd}


\usepackage[%
sorting=nyt, %
backend=biber, %
style=numeric,%
defernumbers=true,%
% sorting=ydnt,%
url=false,%
eprint=false,%
giveninits=true,
maxbibnames=9,%
% doi=false,%
isbn=false,%
sortcites = true,%
citestyle=numeric-comp,
]{biblatex}
\addbibresource{source/moreVoltageOperations.bib}






\begin{document}

\title{Two new chiral 4-polytopes of full rank}
\author{Javier Bracho} 
\address{Institute of Mathematics, National Autonomous University of Mexico (IM UNAM), 04510 Mexico City, Mexico}
\email{jbracho@im.unam.mx}
\author{Daniel Gonz\'alez-Casanova}
\address{Institute of Mathematics, National Autonomous University of Mexico (IM UNAM), 04510 Mexico City, Mexico}
\email{dag@ciencias.unam.mx}
\author{Isabel Hubard}
\address{Institute of Mathematics, National Autonomous University of Mexico (IM UNAM), 04510 Mexico City, Mexico}
\email{isahubard@im.unam.mx}


\section{Introduction}
Polygons and polyhedra are foundational objects in mathematics, yet their formal definitions have evolved significantly. The skeletal approach to polytopes---where structures are viewed as 1-dimensional complexes rather than solid bodies---allows for the emergence of highly symmetric structures not visible in classical definitions. Among these are \emph{chiral polytopes}, which are defined by the presence of maximal rotational symmetry and the absence of reflection symmetry.

While regular polytopes are transitive on flags, chiral polytopes have two flag orbits such that adjacent flags belong to distinct orbits. The study of chiral polytopes in $\mathbb{E}^3$ was pioneered by Schulte (2004--2005), and the first finite example in $\mathbb{E}^4$---Roli's cube---was introduced by Bracho, Hubard, and Pellicer (2014).

Here, we construct two new chiral 4-polytopes in $\mathbb{E}^4$ based on spherical chiral polyhedra identified in \cite{BHP2021}. We demonstrate that the natural rotational symmetries of these polyhedra give rise to distinct chiral 4-polytopes and verify their combinatorial chirality.

\section{Historical Background}
The evolution of polytope theory spans millennia, from the Platonic solids studied in antiquity to the abstract combinatorial definitions of today. Euclid's \emph{Elements} formalized the study of regular polygons and solids. The Renaissance and early modern periods saw the introduction of star polyhedra by Kepler and Poinsot.

In the 19th century, Sch\"afli classified regular polytopes in higher dimensions, and in the 20th century, Coxeter and Gr\"unbaum extended the concept by allowing infinite and non-planar faces. Abstract polytopes, introduced through the study of face lattices and partially ordered sets, became a central framework for understanding these structures, leading to the formal study of chirality by Schulte and Weiss (1991).

\section{Skeletal and Abstract 4-Polytopes in $\mathbb{E}^4$}
A skeletal 4-polytope in $\mathbb{E}^4$ consists of vertices, edges, faces (edge-cycles), and cells (skeletal polyhedra), satisfying conditions of incidence, connectivity, and local finiteness. Each vertex has a well-defined vertex-figure---a connected skeletal polyhedron.

Associated with every skeletal polytope is an abstract polytope, defined as a ranked poset whose elements correspond to faces of various dimensions. A flag is a maximal chain in this poset, and symmetries act by permuting these flags.

\section{Regular and Chiral Abstract 4-Polytopes}
Abstract polytopes are classified as \emph{regular} if their automorphism group acts transitively on flags and as \emph{chiral} if the flag set splits into two orbits with adjacent flags in distinct orbits.

Using string C-groups (for regular polytopes) and their rotational subgroups (for chiral polytopes), we define polytopes via group-theoretic data satisfying specific relations and intersection properties. Faithful geometric realizations in $\mathbb{E}^4$ are constructed via Wythoff's method.

\section{Realization of $\{\frac{5}{2}, 3, 5\}$}
We construct a regular 4-polytope of Schl\"afli type $\{\frac{5}{2}, 3, 5\}$ using a known realization of $\{3,3,5\}$. By identifying a star-shaped icosahedral structure in its vertex-figure and defining appropriate reflections, we show that the associated group satisfies the relations of a string C-group, yielding a realization of $\{\frac{5}{2}, 3, 5\}$.

\section{A Chiral 4-Polytope from $\{\frac{5}{2}, 3, 5\}$}
By composing symmetries of the polytope $\{\frac{5}{2}, 3, 5\}$, we define a subgroup generated by three elements $S_1, S_2, S_3$ satisfying the relations of a chiral string group. The realization obtained via Wythoff's method yields a skeletal chiral 4-polytope, denoted $\{\frac{12_{1,5}}, 3, 5\}$, with helical-faced cells isomorphic to the polyhedron $H_1\left( \left\{ 5, 3, \frac{5}{2} \right\} \right)$. Computational checks confirm its combinatorial chirality.

\section{A Chiral 4-Polytope from $\{5, 3, \frac{5}{2}\}$}
Using the dual structure of the previous polytope, we construct a chiral 4-polytope with cells isomorphic to $H_0\left( \left\{ 5, 3, \frac{5}{2} \right\} \right)$. The abstract group structure and realization process are analogous, and the resulting chiral 4-polytope, denoted $\{\frac{12_{1,5}}, 3, \frac{5}{2}\}$, is shown to be both geometrically and combinatorially chiral.

\section{Computational Verification}
We employ GAP to verify the group relations, intersection properties, and face counts. Realizations are constructed using Wythoff's method based on the coordinates of base flags derived from the polytope $\{3,3,5\}$. Code and data are available at: \href{https://github.com/dan-gc/tesina}{github.com/dan-gc/tesina}.

\section*{Acknowledgments}
The author gratefully acknowledges the guidance and inspiration of Rolf (Roli), Isabel, and Bris; the camaraderie of the \emph{politoperes} group; and the enduring support of Vinicio G\'omez and family.

\begin{thebibliography}{99}
\bibitem{BHP2014} J. Bracho, I. Hubard, D. Pellicer, \emph{A Finite Chiral 4-Polytope in $\mathbb{R}^4$}, Discrete Comput. Geom. 52(4): 799--805 (2014).
\bibitem{BHP2021} J. Bracho, I. Hubard, D. Pellicer, \emph{Chiral Polyhedra in 3-Dimensional Geometries and from a Petrie--Coxeter Construction}, Discrete Comput. Geom. 66(3): 1025--1052 (2021).
% Additional entries as in original thesis...
\end{thebibliography}

\end{document}